\chapter{1989年上海卷}

\section{填空题}

\begin{problem}
若点 \(P(1,a)\) 在曲线 \(x^2+2xy-5y=0\) 上，则 \(a=\)\fillinblank{}.

\end{problem}

\begin{answer}
\(\frac{1}{3}\)
\end{answer}

\begin{problem}
计算 \(2\left( \cos \frac{\pi}{3}+\i \sin \frac{\pi}{3} \right)^5\div \sqrt{2}\left( \cos \frac{\pi}{6}+\i \sin \frac{\pi}{6} \right)=\)\fillinblank{}（\(\i\) 为虚数单位）.

\end{problem}

\begin{answer}
\(-\sqrt{2}\i\)
\end{answer}

\begin{problem}
三角方程 \(\cos\left( \frac{x}{2}+\frac{\pi}{3} \right)=-\frac{1}{2}\) 的解集是 \fillinblank{}.

\end{problem}

\begin{answer}
\(\left\{x\mid x=4k\pi+\frac{2\pi}{3}\text{ 或 }x=4k\pi-2\pi,\ k\in\Z\right\}\)
\end{answer}

\begin{problem}
侧棱长为 \(3\mathrm{cm}\)，底面边长为 \(4\mathrm{cm}\) 的正四棱锥的体积为 \fillinblank{}\(\mathrm{cm}^3\).

\end{problem}

\begin{answer}
\(\frac{16}{3}\)
\end{answer}

\begin{problem}
函数 \(y=x^{-2/3}\) 的递增区间是 \fillinblank{}.

\end{problem}

\begin{answer}
\((-\infty,0)\)
\end{answer}

\begin{problem}
已知长方体 \(ABCD-A'B'C'D'\) 中，棱 \(AA'=5\)，\(AB=12\)，那么直线 \(B'C'\) 和平面 \(A'BCD'\) 的距离是 \fillinblank{}.

\bitmapfigure[max width=0.42\linewidth,max totalheight=4.2cm]{img/1989/shanghai/q6_fig1.png}

\end{problem}

\begin{answer}
\(\frac{60}{13}\)
\end{answer}

\begin{problem}
方程 \(\log_2(x-3)=\log_6(5-x)\) 的解是 \fillinblank{}.

\end{problem}

\begin{answer}
\(x=4\)
\end{answer}

\begin{problem}
已知 \(\sin \alpha=\frac{1}{3}\)，\(2\pi<\alpha<3\pi\)，那么 \(\sin \frac{\alpha}{2}+\cos \frac{\alpha}{2}=\)\fillinblank{}.

\end{problem}

\begin{answer}
\(-\frac{2}{3}\sqrt{3}\)
\end{answer}

\begin{problem}
计算 \(\displaystyle\lim_{n\to\infty}\left( \frac{1}{n^2}+\frac{4}{n^2}+\frac{7}{n^2}+\dots+\frac{3n-2}{n^2} \right)=\)\fillinblank{}.

\end{problem}

\begin{answer}
\(\frac{3}{2}\)
\end{answer}

\begin{problem}
抛物线 \(y=4x^2-4x\) 的准线方程是 \fillinblank{}.

\end{problem}

\begin{answer}
\(y=-\frac{17}{16}\)
\end{answer}

\section{单选题}

\begin{problem}
若 \(\alpha\) 是第四象限的角，则 \(\pi-\alpha\) 是（\quad）

\choices
    {第一象限的角}
    {第二象限的角}
    {第三象限的角}
    {第四象限的角}

\end{problem}

\begin{answer}
C
\end{answer}

\begin{problem}
以点 \(A(-5,4)\) 为圆心，且与 \(x\) 轴相切的圆的标准方程是（\quad）

\choices
    {\((x+5)^2+(y-4)^2=16\)}
    {\((x-5)^2+(y+4)^2=16\)}
    {\((x+5)^2+(y-4)^2=25\)}
    {\((x-5)^2+(y+4)^2=25\)}

\end{problem}

\begin{answer}
A
\end{answer}

\begin{problem}
若 \(a<b<0\)，则下列不等关系中不能成立的是（\quad）

\choices
    {\(\frac{1}{a}>\frac{1}{b}\)}
    {\(\frac{1}{a-b}>\frac{1}{a}\)}
    {\(|a|>|b|\)}
    {\(a^2>b^2\)}

\end{problem}

\begin{answer}
B
\end{answer}

\begin{problem}
两排座位，第一排有 \(3\) 个座位，第二排有 \(5\) 个座位，若 \(8\) 名学生入座（每人一个座位），则不同坐法的种数是（\quad）

\choices
    {\(\mathrm C_8^5 \mathrm C_8^3\)}
    {\(\mathrm A_2^1 \mathrm C_8^5 \mathrm C_8^3\)}
    {\(\mathrm A_8^5 \mathrm A_8^3\)}
    {\(\mathrm A_8^8\)}

\end{problem}

\begin{answer}
D
\end{answer}

\begin{problem}
设 \(z_1\)、\(z_2\) 为复数，那么 \(z_1^2+z_2^2=0\) 是 \(z_1\)、\(z_2\) 同时为零的（\quad）

\choices
    {充分不必要条件}
    {必要不充分条件}
    {充分必要条件}
    {既不充分又不必要条件}

\end{problem}

\begin{answer}
B
\end{answer}

\begin{problem}
下面四个函数中为奇函数的是（\quad）

\choices
    {\(y=x^2\sin\left( x+\frac{\pi}{2} \right)\)}
    {\(y=x^2\cos\left( x+\frac{\pi}{4} \right)\)}
    {\(y=\cos(\operatorname{arc}\operatorname{ctg}x)\)}
    {\(y=\operatorname{arc}\operatorname{ctg}\sin x\)}

\end{problem}

\begin{answer}
C
\end{answer}

\begin{problem}
下列四个命题中的真命题是（\quad）

\choices
    {若直线 \(l\) 与平面 \(\alpha\) 内两条平行直线垂直，则 \(l\perp \alpha\)}
    {若平面 \(\alpha\) 内两条直线与平面 \(\beta\) 内两条直线分别平行，则 \(\alpha\parallel \beta\)}
    {若平面 \(\alpha\) 与直二面角 \(\beta-MN-\gamma\) 的棱 \(MN\) 交于点 \(A\)，与二面角的面 \(\beta\)、面 \(\gamma\) 分别交于 \(AB\)、\(AC\)，则 \(\angle BAC\leqslant 90^\circ\)}
    {以上三个命题都是假命题}

\end{problem}

\begin{answer}
D
\end{answer}

\begin{problem}
要得到函数 \(y=\cos\left( 2x-\frac{\pi}{4} \right)\) 的图象，只需将函数 \(y=\sin 2x\) 的图象（\quad）

\choices
    {向左平移 \(\frac{\pi}{8}\) 个单位}
    {向右平移 \(\frac{\pi}{8}\) 个单位}
    {向左平移 \(\frac{\pi}{4}\) 个单位}
    {向右平移 \(\frac{\pi}{4}\) 个单位}

\end{problem}

\begin{answer}
A
\end{answer}

\begin{problem}
设直角 \(\triangle ABC\) 的直角边 \(BC=a\)，\(AC=b\)，斜边 \(AB=c\)，且 \(a<b\)，现分别以直线 \(BC\)、\(AC\) 和 \(AB\) 为轴将直角三角形绕轴旋转一周，所得的三个旋转体的体积分别记为 \(V_1\)、\(V_2\) 和 \(V_3\)，那么 \(V_1\)、\(V_2\) 和 \(V_3\) 间的大小关系是（\quad）

\choices
    {\(V_3>V_1>V_2\)}
    {\(V_1>V_2>V_3\)}
    {\(V_1>V_3>V_2\)}
    {\(V_2>V_1>V_3\)}

\end{problem}

\begin{answer}
B
\end{answer}

\begin{problem}
下列参数方程（\(t\) 是参数）中与方程 \(y^2=x\) 表示同一曲线的是（\quad）

\choices
    {\(\left\{\begin{aligned}x&=t,\\ y&=t^2.\end{aligned}\right.\)}
    {\(\left\{\begin{aligned}x&=\sin^2 t,\\ y&=\sin t.\end{aligned}\right.\)}
    {\(\left\{\begin{aligned}x&=t,\\ y&=\sqrt{|t|}.\end{aligned}\right.\)}
    {\(\left\{\begin{aligned}x&=\frac{1-\cos 2t}{1+\cos 2t},\\ y&=\operatorname{tg} t.\end{aligned}\right.\)}

\end{problem}

\begin{answer}
D
\end{answer}

\section{解答题}

\begin{problem}
求 \(\cos\left( 2\operatorname{arc}\operatorname{tg}\sqrt{6}+\frac{\pi}{3} \right)\) 的值.

\end{problem}

\begin{solution}
设 \(\theta=\operatorname{arc}\operatorname{tg}\sqrt{6}\)，则 \(\operatorname{tg}\theta=\sqrt{6}\).
\[
\cos 2\theta=\frac{1-\operatorname{tg}^2\theta}{1+\operatorname{tg}^2\theta}=-\frac{5}{7}
\]
\[
\sin 2\theta=\frac{2\operatorname{tg}\theta}{1+\operatorname{tg}^2\theta}=\frac{2\sqrt{6}}{7}
\]
\[
\cos\left( 2\operatorname{arc}\operatorname{tg}\sqrt{6}+\frac{\pi}{3} \right)=\cos 2\theta \cos \frac{\pi}{3}-\sin 2\theta \sin \frac{\pi}{3}=-\frac{5+6\sqrt{2}}{14}
\]
\end{solution}

\begin{problem}
求函数 \(y=\sqrt{2+\log_{1/2}x}+\sqrt{\operatorname{tg}x}\) 的定义域.

\end{problem}

\begin{solution}
解不等式组 \(\left\{\begin{aligned}2+\log_{1/2}x&\geqslant 0,\\ \operatorname{tg}x&\geqslant 0.\end{aligned}\right.\) \circled{1} \circled{2}. 由 \circled{1}，得 \(0<x\leqslant 4\)；由 \circled{2}，得 \(k\pi\leqslant x<k\pi+\frac{\pi}{2}\)，\(k\in\Z\). 所以函数 \(y\) 的定义域为 \(\left( 0,\frac{\pi}{2} \right)\cup [\pi,4]\).
\end{solution}

\begin{problem}
在直角坐标系 \(xOy\) 中，过点 \(P(-3,4)\) 的直线 \(l\) 与直线 \(OP\) 的夹角为 \(45^\circ\)，求直线 \(l\) 的方程.

\end{problem}

\begin{solution}
设直线 \(l\) 的斜率为 \(k\). 直线 \(OP\) 的斜率 \(k_{OP}=-\frac{4}{3}\). 按题意，有
\[
\left| \frac{k+4/3}{1-4k/3} \right|=1
\]
即 \(3k+4=\pm(3-4k)\). 解得 \(k_1=-\frac{1}{7}\)，\(k_2=7\). 故所求直线方程是 \(x+7y-25=0\) 和 \(7x-y+25=0\).
\end{solution}

\begin{problem}
在 \((1-x^2)^{20}\) 的展开式中，如果第 \(4r\) 项和第 \(r+2\) 项的二项式系数相等，（1）求 \(r\) 的值；（2）写出展开式中的第 \(4r\) 项和第 \(r+2\) 项.

\end{problem}

\begin{solution}
（1）展开式中第 \(4r\) 项的二项式系数为 \(\mathrm C_{20}^{4r-1}\)，第 \(r+2\) 项的二项式系数为 \(\mathrm C_{20}^{r+1}\). 根据二项式系数的性质，当且仅当 \(4r-1=r+1\) 或 \((4r-1)+(r+1)=20\) 时，它们的二项式系数相等. 解得 \(r=\frac{2}{3}\)（不合题意，舍去），或 \(r=4\).

（2）当 \(r=4\) 时，第 \(4r\) 项是
\[
T_{16}=\mathrm C_{20}^{15}\left( -x^2 \right)^{15}=-15504x^{30}
\]
第 \(r+2\) 项是
\[
T_6=\mathrm C_{20}^5\left( -x^2 \right)^5=-15504x^{10}
\]
\end{solution}

\begin{problem}
如图，设点 \(A\)、\(B\) 分别在二面角 \(M-PQ-N\)（平面角为锐角）的两个面 \(N\)、\(M\) 上，直线 \(AB\) 与面 \(M\)、\(N\) 所成的角分别为 \(\alpha\)、\(\beta\)，过点 \(A\)、\(B\) 分别作棱 \(PQ\) 的垂线 \(AE\)、\(BF\)，垂足为 \(E\)、\(F\). 求证：\(\frac{AE}{BF}=\frac{\sin \alpha}{\sin \beta}\).

\bitmapfigure[max width=0.34\linewidth,max totalheight=4.8cm]{img/1989/shanghai/q5_fig1.png}

\end{problem}

\begin{solution}
过 \(A\)、\(B\) 分别作 \(AC\)、\(BD\) 垂直于平面 \(M\)、\(N\)，垂足分别为 \(C\)、\(D\)，连 \(CE\)、\(DF\)、\(BC\) 和 \(DA\)，则
\[
\angle ABC=\alpha,\quad \angle BAD=\beta
\]
在 \(\mathrm{Rt}\triangle ABC\) 和 \(\mathrm{Rt}\triangle ABD\) 中，分别有
\[
AC=AB\sin\alpha,\quad BD=AB\sin\beta \tag{\circled{1}}
\]

\bitmapfigure[max width=0.30\linewidth,max totalheight=4.8cm]{img/1989/shanghai/q5_solution_fig1.png}

因为 \(AE\perp PQ\)，所以 \(CE\perp PQ\)（三垂线定理逆定理）. 故 \(\angle AEC\) 为二面角 \(M-PQ-N\) 的平面角，同理 \(\angle BFD\) 也为二面角 \(M-PQ-N\) 的平面角，于是 \(\angle AEC=\angle BFD\)，
\[
\mathrm{Rt}\triangle AEC\sim \mathrm{Rt}\triangle BFD
\]
得 \(\frac{AE}{BF}=\frac{AC}{BD}\). 将 \circled{1} 式代入，即有
\[
\frac{AE}{BF}=\frac{\sin \alpha}{\sin \beta}
\]
\end{solution}

\begin{problem}
设复数集合 \(M=\left\{z\mid |z-2+\i|\leqslant 2,\ z\in \mathbf{C}\right\}\cap \left\{z\mid |z-2-\i|=|z-4+\i|,\ z\in \mathbf{C}\right\}\). （1）试在复平面内作出表示集合 \(M\) 的图形，并说出图形的名称；（2）求集合 \(M\) 中元素 \(z\) 的辐角主值的范围；（3）求集合 \(M\) 中元素 \(z\) 的模的范围.

\end{problem}

\begin{solution}
（1）设 \(z=x+y\i\)，由题意得
\[
\left\{
\begin{aligned}
(x-2)^2+(y+1)^2&\leqslant 4,\\
x-y&=3
\end{aligned}
\right.
\]

\bitmapfigure[max width=0.24\linewidth,max totalheight=4.8cm]{img/1989/shanghai/q6_solution_fig1.png}

则集合 \(M\) 的图形是图中直线段 \(AB\).

（2）解方程组 \(\left\{\begin{aligned}(x-2)^2+(y+1)^2&=4,\\ x-y&=3,\end{aligned}\right.\) 得 \(A(2+\sqrt{2},\sqrt{2}-1)\)，\(B(2-\sqrt{2},-\sqrt{2}-1)\).

记点 \(A\)、\(B\) 所对应的复数的辐角主值分别为 \(\theta_1\)、\(\theta_2\)，则
\[
\theta_1=\operatorname{arc}\operatorname{tg}\frac{3\sqrt{2}-4}{2}
\]
\[
\theta_2=2\pi-\operatorname{arc}\operatorname{tg}\frac{3\sqrt{2}+4}{2}
\]
所以集合 \(M\) 中元素 \(z\) 的辐角主值的范围是
\[
\left[ 2\pi-\operatorname{arc}\operatorname{tg}\frac{3\sqrt{2}+4}{2},2\pi \right)\cup \left[ 0,\operatorname{arc}\operatorname{tg}\frac{3\sqrt{2}-4}{2} \right]
\]

（3）\(|OA|=\sqrt{9+2\sqrt{2}}\)，\(|OB|=\sqrt{9-2\sqrt{2}}\). 作 \(OP\perp AB\)，垂足为 \(P\). 由 \(\left\{\begin{aligned}x-y&=3,\\ x+y&=0\end{aligned}\right.\) 解得 \(P\left( \frac{3}{2},-\frac{3}{2} \right)\)，\(|OP|=\frac{3\sqrt{2}}{2}\).

因为点 \(P\) 在圆 \((x-2)^2+(y+1)^2=4\) 内，又 \(|OA|>|OB|\)，所以集合 \(M\) 中元素 \(z\) 的模的范围是 \(\left[ \frac{3\sqrt{2}}{2},\sqrt{9+2\sqrt{2}} \right]\).
\end{solution}

\begin{problem}
已知等比数列 \(\{a_n\}\) 的首项 \(a_1>0\)，公比 \(q>-1\) 且 \(q\neq 0\). 设数列 \(\{b_n\}\) 的通项 \(b_n=a_{n+1}+a_{n+2}\)（\(n\in N\)），数列 \(\{a_n\}\)、\(\{b_n\}\) 的前 \(n\) 项的和分别记为 \(A_n\)、\(B_n\). 试比较 \(A_n\) 与 \(B_n\) 的大小.

\end{problem}

\begin{solution}
解法一：因为 \(b_n=a_{n+1}+a_{n+2}=a_n(q^2+q)\)，
\[
B_n=(q^2+q)A_n
\]
当 \(q>0\) 时，\(a_n=a_1q^{n-1}>0\)，所以 \(A_n>0\)；当 \(-1<q<0\) 时，
\[
A_n=\frac{a_1(1-q^n)}{1-q}>0
\]
故当 \(q>-1\) 且 \(q\neq 0\) 时总有 \(A_n>0\). 由
\[
B_n-A_n=A_n(q^2+q-1)=A_n\left( q+\frac{\sqrt{5}+1}{2} \right)\left( q-\frac{\sqrt{5}-1}{2} \right)
\]
得：

（1）当 \(-1<q<\frac{\sqrt{5}-1}{2}\) 且 \(q\neq 0\) 时，\(A_n>B_n\)；

（2）当 \(q=\frac{\sqrt{5}-1}{2}\) 时，\(A_n=B_n\)；

（3）当 \(q>\frac{\sqrt{5}-1}{2}\) 时，\(A_n<B_n\).

解法二：因为 \(b_n=a_{n+1}+a_{n+2}=a_1(1+q)q^n\)，
\[
\frac{b_{n+1}}{b_n}=q,\quad q\neq 0,\quad n\in N
\]
因此数列 \(\{b_n\}\) 也是等比数列.

\circled{1} 当 \(q=1\) 时，
\[
A_n=na_1,\quad B_n=2na_1
\]
所以 \(A_n<B_n\).

\circled{2} 当 \(q>-1\) 且 \(q\neq 1\) 时，
\[
B_n-A_n=\frac{a_1(1-q^n)(q^2+q-1)}{1-q}
=\frac{a_1(1-q^n)\left( q+\frac{\sqrt{5}+1}{2} \right)\left( q-\frac{\sqrt{5}-1}{2} \right)}{1-q}
\]
因为在 \(-1<q<1\) 时，\(1-q^n>0\)，\(1-q>0\)；在 \(q>1\) 时，\(1-q^n<0\)，\(1-q<0\). 又 \(a_1>0\)，所以 \(\frac{a_1(1-q^n)}{1-q}>0\). 于是：

（1）当 \(-1<q<\frac{\sqrt{5}-1}{2}\) 且 \(q\neq 0\) 时，\(A_n>B_n\)；

（2）当 \(q=\frac{\sqrt{5}-1}{2}\) 时，\(A_n=B_n\)；

（3）当 \(q>\frac{\sqrt{5}-1}{2}\) 时，\(A_n<B_n\).
\end{solution}

\begin{problem}
\(F\) 是定点，\(l\) 是定直线，点 \(F\) 到直线 \(l\) 的距离为 \(p\)（\(p>0\)），点 \(M\) 在直线 \(l\) 上滑动，动点 \(N\) 在 \(MF\) 的延长线上，且满足条件 \(\frac{|FN|}{|MN|}=\frac{1}{|MF|}\). （1）求动点 \(N\) 的轨迹；（2）求 \(|MN|\) 的最小值.

\end{problem}

\begin{solution}
（1）解法一：作 \(FK\perp l\)，垂足为 \(K\)，以 \(F\) 为极点，\(FK\) 的反向延长线为极轴，建立极坐标系. 设动点 \(N\) 的极坐标为 \((\rho,\theta)\). 根据题意，不妨取 \(\rho>0\)，\(\cos\theta>0\).

\bitmapfigure[max width=0.22\linewidth,max totalheight=4.8cm]{img/1989/shanghai/q7_solution_fig1.png}

则 \(|MF|=\frac{p}{\cos\theta}\)，\(|FN|=\rho\)，从而
\[
|MN|=|MF|+|FN|=\frac{p}{\cos\theta}+\rho
\]
由点 \(N\) 所满足的条件，得
\[
\frac{p}{\cos\theta}\cdot \rho=\frac{p}{\cos\theta}+\rho
\]
故所求轨迹的极坐标方程为
\[
\rho=\frac{1}{1-\frac{1}{p}\cos\theta}\quad (0<\cos\theta<p)
\]
设过极点 \(F\) 且与极轴垂直的直线为 \(l'\)，则当 \(e=\frac{1}{p}>1\)，即 \(0<p<1\) 时，所求轨迹是双曲线在直线 \(l'\) 右边的部分；当 \(e=\frac{1}{p}=1\)，即 \(p=1\) 时，所求轨迹是抛物线在直线 \(l'\) 右边的部分；当 \(0<e=\frac{1}{p}<1\)，即 \(p>1\) 时，所求轨迹是椭圆在直线 \(l'\) 右边的部分.

解法二：过 \(F\) 作 \(l\) 的垂线，垂足为 \(K\)，以 \(F\) 为原点，直线 \(KF\) 为 \(x\) 轴，建立直角坐标系. 设点 \(N\) 的坐标为 \((x,y)\)，根据题意，\(x>0\). 直线 \(l\) 的方程为 \(x=-p\). 设直线 \(MN\) 的斜率为 \(k\)，则点 \(M\) 和 \(N\) 的坐标分别为 \((-p,-pk)\) 和 \((x,kx)\).

\bitmapfigure[max width=0.22\linewidth,max totalheight=4.8cm]{img/1989/shanghai/q7_solution_fig2.png}

有 \(|MF|=p\sqrt{1+k^2}\)，\(|FN|=\sqrt{x^2+y^2}=x\sqrt{1+k^2}\)，
\[
|MN|=|MF|+|FN|=(x+p)\sqrt{1+k^2}
\]
按题意有 \((p\sqrt{1+k^2})\cdot (x\sqrt{1+k^2})=(x+p)\sqrt{1+k^2}\). 因为 \(k=\frac{y}{x}\)，所以
\[
p\sqrt{x^2+y^2}=x+p \tag{\circled{1}}
\]
化简得
\[
(p^2-1)x^2+p^2y^2-2px-p^2=0\quad (x>0)
\]
当 \(0<p<1\) 时，所求轨迹是双曲线在 \(y\) 轴右边的部分；当 \(p=1\) 时，所求轨迹是抛物线在 \(y\) 轴右边的部分；当 \(p>1\) 时，所求轨迹是椭圆在 \(y\) 轴右边的部分.

（2）解法一：
\[
|MN|=\frac{p}{\cos\theta}+\frac{1}{1-\frac{1}{p}\cos\theta}=\frac{p^2}{\cos\theta (p-\cos\theta)}=\frac{p^2}{-\left( \cos\theta-\frac{p}{2} \right)^2+\frac{p^2}{4}}
\]
若 \(0<p\leqslant 2\)，则当 \(\cos\theta=\frac{p}{2}\) 时，\(|MN|\) 取得最小值 \(4\)；若 \(p>2\)，则当 \(\cos\theta=1\) 时，\(|MN|\) 取得最小值 \(\frac{p^2}{p-1}\).

解法二：将 \(y=kx\) 代入 \circled{1}，得
\[
x=\frac{p}{p\sqrt{1+k^2}-1}\quad \text{（由 \(x>0\)，得 \(k^2>\frac{1-p^2}{p^2}\)）}
\]
\[
|MN|=(x+p)\sqrt{1+k^2}=\frac{p^2(1+k^2)}{p\sqrt{1+k^2}-1}
\]
令 \(k=\operatorname{tg}\theta\left( -\frac{\pi}{2}<\theta<\frac{\pi}{2} \right)\)，由 \(\operatorname{tg}^2\theta>\frac{1-p^2}{p^2}\)，得 \(0<\cos\theta<p\). 则
\[
|MN|=\frac{p^2}{\cos\theta (p-\cos\theta)}
\]
以下同解法一.
\end{solution}
